\documentclass[10pt,a4paper]{amsart}
\usepackage[margin=19mm]{geometry}
\usepackage{amsmath,amssymb,mathtools}
\usepackage[scr=rsfs,cal=euler]{mathalfa}
\usepackage{xfrac}
\usepackage[unicode,hidelinks,bookmarks=false]{hyperref}
\newcommand{\ie}{i.e.\ }
\newcommand{\cf}{cf.\ }
\newcommand{\resp}{respectively\ }
\newcommand{\Subsection}{\subsection}
\providecommand{\nobreakdash}{\nobreak\hbox{-}}
\setlength{\emergencystretch}{2em}
\multlinegap0pt
\usepackage{bm}
\usepackage{bbm}
\usepackage{dsfont}
\usepackage{xspace}
\usepackage{cancel}
\usepackage{mathtools}
\usepackage{amsthm}
\makeatletter
\@ifundefined{thm}{\newtheorem{thm}{Thm}}{}
\@ifundefined{assum}{\newtheorem{assum}[thm]{Assumptions}}{}
\@ifundefined{prop}{\newtheorem{prop}[thm]{Proposition}}{}
\makeatother
\title[Existence for the Discrete Nonlinear Fragmentation Equation ]{Existence for the Discrete Nonlinear Fragmentation Equation with Degenerate Diffusion}
\author{Saumyajit Das, Ram Gopal Jaiswal}
\date{Source: arXiv:2602.14070v3 (CC BY 4.0)}
\begin{document}
\maketitle

\noindent\emph{Source and licence.} Existence for the Discrete Nonlinear Fragmentation Equation with Degenerate Diffusion. Version arXiv:2602.14070v3 (CC BY 4.0), distributed under the Creative Commons Attribution 4.0 International licence, as recorded in the arXiv licence metadata for this exact version and shown on the abstract page https://arxiv.org/abs/2602.14070v3. Full rights notice: licenses/arxiv-2602.14070-CC-BY-4.0.txt.\par\smallskip

\noindent\textit{Editorial scope.} Counted unit: the Prop whose source label is \texttt{epsilon compactness in the inequation} in the article (source0.tex lines 2570--2601, source label \texttt{epsilon compactness in the inequation}), delivered together with the article's own complete original proof of it (source0.tex lines 2602--2936, one \texttt{proof} environment of 335 lines). No mathematics is written, repaired or re-derived: statement and proof are the article's own text line for line. Required context retained uncounted: NFE Operator (lines 557--564); strong assumption (lines 693--725); eq:RNFE (lines 1104--1117); prop:dual\_estimate (lines 1575--1596); m\_trunction (lines 1744--1748); representation regularization (lines 2199--2202); gradient estimate:regularization (lines 2290--2299). Every definition, display and label of those lines is retained unchanged. Omitted from this package and not redistributed: 1--556, 565--692, 726--1103, 1118--1574, 1597--1743, 1749--2198, 2203--2289, 2300--2569, 2937--3386. Nothing inside a retained excerpt is omitted or altered: every retained body is delivered exactly as the article's lines are written, so no omission receipt of any kind accompanies this package. Citations: the retained excerpts carry no citation command at all, so no bibliography entry of the article is transcribed and no reference is redistributed. Reference adaptation: none was needed and none was made; every reference inside the delivered unit resolves inside it, so no unresolved reference or citation is produced. Printed slips kept literal: no printed word, symbol, hypothesis or number of the counted statement or of its proof was altered. Layout and class: the article's own class is \texttt{amsart} with packages including amsmath, amsfonts, amssymb, amsthm, enumerate, multicol, mathrsfs, tikz, graphicx, hyperref, caption, enumitem, wasysym, cleveref; the wrapper is the lane's pinned \texttt{amsart} editorial wrapper at 10pt on A4, which restates the theorem-like environments the retained text needs and adds the article's own macro definitions that the retained text needs (none), with extra wrapper packages (bm, bbm, dsfont, xspace, cancel, mathtools). The delivered numbering is the unit's own. Assets: no retained span contains a figure environment, a graphics-inclusion command, a PDF-page inclusion, a table or a TikZ picture; no image file is copied into this package, the package declares no asset dependency and no image file is redistributed with it. Licence: version arXiv:2602.14070v3 of this article is distributed under the Creative Commons Attribution 4.0 International licence, as recorded in the arXiv licence metadata for this exact version and shown on the abstract page https://arxiv.org/abs/2602.14070v3. A full rights notice is delivered at licenses/arxiv-2602.14070-CC-BY-4.0.txt. Characters: every retained excerpt is pure ASCII, so the delivered engine, XeLaTeX with its default Latin Modern text font, reports no missing character.
	\begin{align} 
		Q_i(f)
		= 
		\frac12 \sum_{j=i+1}^{\infty}\sum_{k=1}^{j-1}
		b_{j-k,k}^i\, a_{j-k,k}\, f_{j-k} f_k
		\;-\;
		\sum_{j=1}^{\infty} a_{i,j}\, f_i f_j. \label{NFE Operator}
	\end{align}

	\begin{assum}\label{strong assumption}
		The following assumptions apply to the collision kernel $a_{i,j}$, the
		breakage function $b_{i,j}^k$, and the diffusion coefficients $d_i$
		
		\begin{enumerate}
			\item[(A1)] The collision kernel $a_{i,j}=a_{j,i}\ge0$ satisfies the summability condition
			\[
			\sum_{i,j \ge 1} (i+j)\, a_{i,j} < \infty.
			\]
			
			\item[(A2)] There exists a constant $B > 0$ such that
			\[
			b_{i,j}^k \le B, \qquad 1 \le k \le i+j-1, \quad i,j \ge 1.
			\]
			
			\item[(A3)] The diffusion coefficients satisfy
			\[
			\lim_{i \to \infty} d_i = 0.
			\]
			
			\item[(A4)] The following summability condition holds
			\[
			\sum_{j=1}^{\infty}
			\sum_{k=1}^{\infty}
			\sum_{i=1}^{j+k-1}
			\frac{\sqrt{\, b^{\,i}_{\,j,k}\, a_{\,j,k}}}{\sqrt{k j\, d_i d_j}}
			\;+\;
			\sum_{i=1}^{\infty}\sum_{j=1}^{\infty}
			\frac{\sqrt{a_{i j}}}{\sqrt{i j\, d_i d_j}}
			< \infty.
			\]
		\end{enumerate}
	\end{assum}

		\begin{equation}
			\label{eq:RNFE}
			\left\{
			\begin{aligned}
				& \partial_t f_{i,\varepsilon} - d_i \Delta f_{i,\varepsilon}
				= Q_{i,\varepsilon}(f_{\varepsilon})
				&& \text{in } (0,T)\times\Omega \\
				& \nabla f_{i,\varepsilon} \cdot \nu = 0
				&& \text{on } (0,T)\times\partial\Omega \\
				& f_{i,\varepsilon}^n(0,x) = f_i^{\mathrm{in}}(x)
				&& \text{in}\ \Omega,
			\end{aligned}
			\right.
		\end{equation}

	\begin{prop}
		\label{prop:dual_estimate}
		Assume that
		\[
		\sum_{i=1}^{\infty} i f_i^{\mathrm{in}} \in L^2(\Omega),
		\qquad
		\sup_{i\in\mathbb{N}^*} d_i < +\infty.
		\]
		Then, for all $T>0$, the solutions to the regularized system \eqref{eq:RNFE} satisfy the following estimate
		\begin{equation}
			\label{eq:dual_estimate}
			\int_0^T \int_\Omega
			\left( \sum_{i=1}^{\infty} i d_i f_{i,\varepsilon}(t,x) \right)
			\left( \sum_{i=1}^{\infty} i f_{i,\varepsilon}(t,x) \right)
			\,dx\,dt
			\;\lesssim\;
			\left( \sup_{i\ge1} d_i \right)
			\left\|
			\sum_{i=1}^{\infty} i f_i^{\mathrm{in}}
			\right\|_{L^2(\Omega)}^{2}.
		\end{equation}
	\end{prop}

	\begin{align}\label{m_trunction}
		T_m(\sigma)=\sigma \quad \text{for } 0\le \sigma \le m-1,
		\qquad
		T_m(\sigma)=m \quad \text{for } \sigma \ge m,
	\end{align}

		\begin{align}\label{representation regularization}
			f_{i,\varepsilon}(t,x)= \int_{\Omega}G_{d_i}(t,0,x,y)f_i^{\rm{in}}(y)\, dy+\int_{0}^{t}\int_{\Omega} G_{d
				_i}(t,s,x,y)Q_{i,\varepsilon}(s,y) \, ds\, dy, \ \ \forall \, i\in\mathbb{N},
		\end{align}

		\begin{align}\label{gradient estimate:regularization}
			\|\nabla f_{i,\varepsilon}\|_{L^1((0,T)\times\Omega)}
			\le
			\frac{C\sqrt{T}}{\sqrt{d_i}}
			\left(
			\|f_i^{\mathrm{in}}\|_{L^1(\Omega)}
			+
			\|Q_{i,\varepsilon}\|_{L^1((0,T)\times\Omega)}
			\right).
		\end{align}

	\begin{prop}\label{epsilon compactness in the inequation}
		Let $m>0$, and let $T_m$ be a smooth function as defined in \eqref{m_trunction}. Let $f_i$ 
		denote the $L^1((0,T); W^{1,1}(\Omega))$ limit of the sequence $\{f_{i,\varepsilon}\}$, 
		where $f_{i,\varepsilon}$ is a solution to the regularized nonlinear fragmentation system~\eqref{eq:RNFE} 
		for all $i \in \mathbb{N}$ and for all $\varepsilon > 0$. Further assume that
		\[
		\sum_{i=1}^{\infty} i f_i^{\mathrm{in}} \in L^2(\Omega),
		\ \text{and} \ 
		\sup_{i\in\mathbb{N}^*} d_i < +\infty.
		\]
		Then, for every nonnegative test function 
		$\psi \in C_c^\infty([0,T)\times\Omega)$, we have
		\begin{align}\label{control estimate on epsilon}
			\int_{\Omega} \psi(0,x)\,T_m(\omega_i^{\mathrm{in}})\,dx
			+&
			\int_{0}^{T}\!\!\int_{\Omega} (-\partial_t\psi)\,T_m(\omega_i)\,dx\,dt
			+
			d_i \int_{0}^{T}\!\!\int_{\Omega} \nabla\psi \, T_m'(\omega_i) \, \nabla\omega_i \,dx\,dt 
			\nonumber  \\
			&\ge
			\int_{0}^{T}\!\!\int_{\Omega} \psi\,\Lambda^m_i\,dx\,dt
			-
			C\,m^2(1+T)^2\,\eta^{1/2},
		\end{align}
		where
		\[
		\omega_i = f_i + \eta \sum_{j\ne i} f_j, 
		\ \text{and} \ 
		\Lambda^m_i = T_m'(\omega_i) \Bigl( Q_i + \eta \sum_{j\ne i} Q_j \Bigr).
		\]
		Here $Q_i$ is as defined in \eqref{NFE Operator}.
	\end{prop}

	\begin{proof}
		We pass the limit $\varepsilon\to 0$ term by term.
		\newline
		\textbf{Convergence of the truncated densities.}
		We consider
		\[
		\int_{0}^{T}\!\!\int_{\Omega}
		\left|
		f_{i,\varepsilon}
		+ \eta \sum_{j\ne i} f_{j,\varepsilon}
		-
		f_i
		- \eta \sum_{j\ne i} f_j
		\right| dx\,dt .
		\]
		Using the triangle inequality, we obtain
		\begin{align}\label{intermediate 1}
			\le
			\int_{0}^{T}\!\!\int_{\Omega}
			|f_{i,\varepsilon}-f_i|\,dx\,dt
			+
			\eta \sum_{j\ne i}
			\int_{0}^{T}\!\!\int_{\Omega}
			|f_{j,\varepsilon}-f_j|\,dx\,dt .
		\end{align}
		Fix $M\in\mathbb{N}$. Recall that the solution $f_{i,\varepsilon}$ can be expressed as \eqref{representation regularization}:
		\[
		f_{i,\varepsilon}(t,x)= \int_{\Omega}G(t,0,x,y)f_i^{\rm{in}}(y)\, dy+\int_{0}^{t}\int_{\Omega} G(t,s,x,y)Q_{i,\varepsilon}(s,y) \, ds\, dy, \ \ \forall \, i\in\mathbb{N}.
		\]
		Hence, we obtain the following $L^1((0,T)\times\Omega)$ integral estimate
		\begin{align*}
			\sum_{i>M} \|f_{i,\varepsilon}(t,x)\|_{L^1((0,T)\times\Omega)} \leq &\sum_{i>M} \int_0^T\| f_i^{\rm{in}}\|_{L^1(\Omega)}+ \sum_{i>M} \left\| \int_{0}^{t}\int_{\Omega} G(t,s,x,y)Q_{i,\varepsilon}(s,y) \, ds\, dy\right\|_{L^1((0,T)\times\Omega)}\\
			\leq &T\sum_{i>M} \| f_i^{\rm{in}}\|_{L^1(\Omega)}+\sum_{i>M} \| Q_{i,\varepsilon}\|_{L^1((0,T)\times\Omega)} \\\leq & T\sum_{i>M} \| f_i^{\rm{in}}\|_{L^1(\Omega)}+C\sum_{i>M}
			\sum_{j=i+1}^{\infty} \sum_{k=1}^{j-1}
			\frac{b_{j-k,k}^i a_{j-k,k}}{k(j-k)d_{k}d_i}+C\sum_{i>M}
			\sum_{j=1}^{\infty}
			\frac{a_{i,j}}{ijd_jd_i} \\
			\leq & \delta_M\to 0 \ \mbox{as}\ M\to \infty.
		\end{align*}
		Here the third line is due to Proposition \ref{prop:dual_estimate}. Since $f_{i,\varepsilon}\ge0$ and $f_{i,\varepsilon}\to f_i$ pointwise a.e.,
		Fatou's lemma yields
		\[
		\sum_{i>M} \|f_{i}\|_{L^1((0,T)\times\Omega)} \leq \sum_{i>M} \|f_{i,\varepsilon}\|_{L^1((0,T)\times\Omega)}\leq \delta_M\to 0 \ \mbox{as}\ M\to \infty.
		\]
		Splitting the sum in \eqref{intermediate 1} into a finite and an infinite part,
		we obtain
		\[
		\le
		\int_{0}^{T}\!\!\int_{\Omega}
		|f_{i,\varepsilon}-f_i|\,dx\,dt
		+
		\eta \sum_{j\ne i}^{M}
		\int_{0}^{T}\!\!\int_{\Omega}
		|f_{j,\varepsilon}-f_j|\,dx\,dt
		+
		\eta\,\delta_M,
		\]
		where $\delta_M\to0$ as $M\to\infty$, uniformly in $\varepsilon$, thanks to
		the uniform $L^1$ summability of $\{f_{j,\varepsilon}\}_{j\ge1}$.
		Since $f_{j,\varepsilon}\to f_j$ strongly in $L^1((0,T)\times\Omega)$ for
		each fixed $j$, we conclude that
		\[
		f_{i,\varepsilon}
		+ \eta \sum_{j\ne i} f_{j,\varepsilon}
		\;\longrightarrow\;
		f_i + \eta \sum_{j\ne i} f_j
		\quad \text{in } L^1((0,T)\times\Omega).
		\]
		In particular,
		\[
		\omega_{i,\varepsilon}
		\longrightarrow
		\omega_i
		\quad \text{a.e. and in } L^1((0,T)\times\Omega).
		\]
		
		\medskip
		\noindent\textbf{Limit of the second term (time derivative).}
		Since $T_m$ is continuous and bounded, we have
		\[
		T_m(\omega_{i,\varepsilon}) \to T_m(\omega_i)
		\quad \text{a.e. in } (0,T)\times\Omega,
		\qquad
		|T_m(\omega_{i,\varepsilon})|\le m .
		\]
		Therefore, by the Dominated Convergence Theorem,
		\[
		\lim_{\varepsilon\to0}
		\int_{0}^{T}\!\!\int_{\Omega}
		(-\partial_t\psi)\,T_m(\omega_{i,\varepsilon})\,dx\,dt
		=
		\int_{0}^{T}\!\!\int_{\Omega}
		(-\partial_t\psi)\,T_m(\omega_i)\,dx\,dt .
		\]
		
		\medskip
		\noindent\textbf{Gradient tail estimation.}
		We recall the following a priori estimates on the gradient term derived earlier \eqref{gradient estimate:regularization}, 
		\[
		\|\nabla f_{i,\varepsilon}\|_{L^1((0,T)\times\Omega)}
		\le
		C\sqrt{T}\,\frac{1}{\sqrt{d_i}}
		\left(
		\|f_i^{\mathrm{in}}\|_{L^1(\Omega)}
		+
		\|Q_{i,\varepsilon}\|_{L^1((0,T)\times\Omega)}
		\right).
		\]
		The sum of the tail can be estimated as 
		\begin{align}\label{gradient uniform sum}
			\sum_{i>M}
			\|\nabla f_{i,\varepsilon}\|_{L^1((0,T)\times\Omega)}
			\le C\sqrt{T}\, \sum_{i>M}\frac{1}{\sqrt{d_i}}
			\left(
			\|f_i^{\mathrm{in}}\|_{L^1(\Omega)}
			+
			\sum_{i>M}\|Q_{i,\varepsilon}\|_{L^1((0,T)\times\Omega)}
			\right)
		\end{align}
		Therefore,
		\begin{align*}
			\left\|
			\frac{Q_{i,\varepsilon}(f)}{d_i}
			\right\|_{L^1((0,T)\times\Omega)}
			&\le
			\frac{1}{2}\left\|
			\sum_{j=i+1}^{\infty} \sum_{k=1}^{j-1}
			\frac{b_{j-k,k}^i a_{j-k,k}}{k(j-k)d_{k}d_i}\,
			(j-k)f_{j-k,\varepsilon} kd_kf_{k,\varepsilon}
			\right\|_{L^1((0,T)\times\Omega)}
			\\
			&\quad+
			\left\|
			\sum_{j=1}^{\infty}
			\frac{a_{i,j}}{ijd_jd_i}\, if_{i,\varepsilon} jd_jf_{j,\varepsilon}
			\right\|_{L^1((0,T)\times\Omega)}.\\
			\lesssim & \frac{1}{2}
			\sum_{j=i+1}^{\infty} \sum_{k=1}^{j-1}
			\frac{b_{j-k,k}^i a_{j-k,k}}{k(j-k)d_{k}d_i}\,
			+ \sum_{j=1}^{\infty}
			\frac{a_{i,j}}{ijd_jd_i},
		\end{align*}
		where we obtain the last inequality thanks to Proposition~\ref{prop:dual_estimate}. Hence the tail 
		\begin{align}\label{tail control source}
			{\sum_{i>M}\left\|
				\frac{Q_{i,\varepsilon}(f)}{d_i}
				\right\|_{L^1((0,T)\times\Omega)}\le \delta_M \to 0}
		\end{align}
		as $M\to\infty$. As a consequence, for $\delta_M>0$, there exists $M\in\mathbb{N}$ such that
		\[
		\sum_{j>M}
		\|\nabla f_{j,\varepsilon}\|_{L^1((0,T)\times\Omega)}
		<\delta_M \to 0 \ \ \text{as}\ M\to+\infty,
		\ \text{uniformly in } \varepsilon.
		\]
		Since $\nabla f_{i,\varepsilon}\to \nabla f_{i}$ pointwise. Fatou's lemma yields
		\[
		\sum_{j>M}
		\|\nabla f_{j}\|_{L^1((0,T)\times\Omega)}
		<\delta_M \to 0 \ \ \text{as}\ M\to+\infty.
		\]
		\medskip
		\noindent\textbf{Convergence of the gradients of $\omega_{i,\varepsilon}$.}
		We consider
		\[
		\int_{0}^{T}\!\!\int_{\Omega}
		\left|
		\nabla f_{i,\varepsilon}
		+ \eta \sum_{j\ne i} \nabla f_{j,\varepsilon}
		-
		\nabla f_i
		-
		\eta \sum_{j\ne i} \nabla f_j
		\right| dx\,dt .
		\]
		Using the triangle inequality and splitting the infinite sum at $M$, we obtain
		\begin{align*}
			\le{}&
			\int_{0}^{T}\!\!\int_{\Omega}
			|\nabla f_{i,\varepsilon}-\nabla f_i|\,dx\,dt
			+
			\eta \sum_{\substack{j=1\\ j\ne i}}^{M}
			\int_{0}^{T}\!\!\int_{\Omega}
			|\nabla f_{j,\varepsilon}-\nabla f_j|\,dx\,dt
			\\
			&\qquad
			+
			\eta \sum_{j>M}
			\left(
			\|\nabla f_{j,\varepsilon}\|_{L^1}
			+
			\|\nabla f_j\|_{L^1}
			\right).
		\end{align*}
		For each fixed $j$, we have $\nabla f_{j,\varepsilon}\to\nabla f_j$ in
		$L^1((0,T)\times\Omega)$, while the tail is controlled uniformly by the
		previous estimate. Hence,
		\[
		\nabla\omega_{i,\varepsilon}
		\longrightarrow
		\nabla\omega_i
		\quad \text{in } L^1((0,T)\times\Omega).
		\]
		
		\medskip
		\noindent\textbf{Limit of the third term (diffusion).}
		Since $T_m'$ is bounded and continuous, it follows that
		\[
		T_m'(\omega_{i,\varepsilon})\,\nabla\omega_{i,\varepsilon}
		\longrightarrow
		T_m'(\omega_i)\,\nabla\omega_i
		\quad \text{in } L^1((0,T)\times\Omega).
		\]
		Therefore,
		\[
		\lim_{\varepsilon\to0}
		d_i \int_{0}^{T}\!\!\int_{\Omega}
		\nabla\psi\,
		T_m'(\omega_{i,\varepsilon})\,
		\nabla\omega_{i,\varepsilon}\,dx\,dt
		=
		d_i \int_{0}^{T}\!\!\int_{\Omega}
		\nabla\psi\,
		T_m'(\omega_i)\,
		\nabla\omega_i\,dx\,dt .
		\]
		
		\medskip
		\noindent\textbf{Limit of the fourth term (reaction).}
		Recall that
		\[
		\Lambda^m_{i,\varepsilon}
		=
		T_m'(\omega_{i,\varepsilon})
		\left(
		Q_{i,\varepsilon}
		+
		\eta \sum_{j\ne i} Q_{j,\varepsilon}
		\right).
		\]
		Fix $(t,x)\in \text{Supp}\, \{ T'_m{(w_{j,\varepsilon})}\}$.  
		By the construction of the truncation $T_m$ \eqref{m_trunction}, we have, for all $j\ge1$ and
		$\varepsilon>0$,
		\[
		0 \le f_{j,\varepsilon}(t,x) \le m .
		\]
		Consequently, 
		\[
		0 \le \varepsilon\sum_{j=1}^{\infty} c_j (f_{j,\varepsilon}(t,x))^2
		\le
		\varepsilon m^2 \sum_{j=1}^{\infty} c_j, \quad \forall \, (t,x)\in \text{Supp}\, \{ T'_m{(w_{j,\varepsilon})}\}. 
		\]
		By the assumption on the coefficients (see assumption~\ref{strong assumption}), we have
		\[
		\sum_{j=1}^{\infty} c_j < \infty .
		\]
		Hence,
		\[
		\varepsilon m^2 \sum_{j=1}^{\infty} c_j
		\;\xrightarrow[\varepsilon\to0]{}\; 0 .
		\]
		The above estimate is pointwise and uniform with respect to $(t,x)\in \text{Supp}\, \{ T'_m{(w_{j,\varepsilon})}\}$.
		As a consequence,
		\[
		1+\varepsilon\sum_{j=1}^{\infty} c_j (f_{j,\varepsilon})^2(t,x)
		\xrightarrow[\varepsilon\to0]{}
		1,
		\quad \forall \, (t,x)\in \text{Supp}\, \{ T'_m{(w_{j,\varepsilon})}\} .
		\]
		Similarly, with the help of uniform tail control \eqref{tail control source} on $Q_{j,\varepsilon}$, for any $j\ge 1$, we obtain
		\[
		T_{m}^{'}(w_{j,\varepsilon})Q_{j,\varepsilon} \to T_m^{'}(w_j)Q_j
		\]
		It yields
		\[
		\Lambda^m_{i,\varepsilon} \to \Lambda^m_i
		:=
		T_m'(\omega_i)
		\left(
		Q_i + \eta \sum_{j\ne i} Q_j
		\right).
		\]
		Moreover, the truncation yields the domination
		\[
		|\psi\,\Lambda^m_{i,\varepsilon}|
		\le
		\left(\frac{m}{\eta}\right)^2
		\left(
		\sum_{i=1}^{\infty}\sum_{j=i+1}^{\infty}\sum_{k=1}^{j-1}
		b^i_{j-k,k}\,a_{j-k,k}
		+
		\sum_{i=1}^{\infty}\sum_{j=1}^{\infty} a_{ij}
		\right),
		\]
		which is integrable. Hence, by the Dominated Convergence Theorem,
		\[
		\Lambda^m_{i,\varepsilon} \to \Lambda^m_i
		\quad \text{in } L^1((0,T)\times\Omega).
		\]
		
		\medskip
		\noindent\textbf{Limit weak formulation.}
		Letting $\varepsilon\to0$ in the weak formulation, we conclude that for every
		nonnegative test function $\psi\in C_c^\infty([0,T)\times\Omega)$,
		\begin{align*}
			\int_{\Omega} \psi(0,x)\,T_m(\omega_i^{\mathrm{in}})\,dx
			+&
			\int_{0}^{T}\!\!\int_{\Omega}
			(-\partial_t\psi)\,T_m(\omega_i)\,dx\,dt
			+
			d_i \int_{0}^{T}\!\!\int_{\Omega}
			\nabla\psi\,
			T_m'(\omega_i)\,
			\nabla\omega_i\,dx\,dt
			\\
			&\qquad \ge
			\int_{0}^{T}\!\!\int_{\Omega}
			\psi\,\Lambda^m_i\,dx\,dt
			-
			C\,m\,\eta^{1/2},
		\end{align*}
		where
		\[
		\omega_i
		=
		f_i + \eta \sum_{j\ne i} f_j,
		\ \text{and} \
		\Lambda^m_i
		=
		T_m'(\omega_i)
		\left(
		Q_i + \eta \sum_{j\ne i} Q_j
		\right).
		\]
	\end{proof}

\end{document}
